\documentclass{article}

\usepackage[margin=1in]{geometry}
\usepackage{booktabs}
\usepackage{array}
\usepackage{xcolor}
\usepackage{hyperref}
\usepackage{amsmath}

\title{Supplementary Material: Spectrogram-Only Self-Supervised Pretraining for Low-Resource Speech Recognition}
\author{Anonymous Author}
\date{}

\newcommand{\todo}[1]{\textcolor{red}{TODO: #1}}
\newcommand{\code}[1]{\texttt{#1}}

\begin{document}
\maketitle

\section{Full Training Procedure}

\subsection{Target Generation}
The unlabeled corpus is LibriSpeech train-960, implemented as the
\code{train-clean-100}, \code{train-clean-360}, and \code{train-other-500}
splits in \code{slurm/causal\_specunit/01\_generate\_targets.sh}. Audio is
loaded at 16 kHz, converted to mono, peak-normalized, pre-emphasized with
coefficient 0.97, and transformed into 80-bin log-mel features. Global
CMVN statistics are estimated over the unlabeled training corpus and then
applied before chunking.

For iter-1, normalized log-mel frames are grouped into chunks of size 8 and
stride 4. A PCA projection with whitening maps chunk vectors to 64
dimensions. MiniBatch $K$-means is then fit twice on up to 1,000,000
sampled chunks, producing a coarse codebook with $K_c=100$ and a fine
codebook with $K_f=500$. The assigned cluster IDs are cached per utterance
and sharded for training.

For iter-2, \code{slurm/causal\_specunit/04\_iter2\_pretrain\_chain.sh}
loads the 150k-step iter-1 checkpoint, runs the encoder over the same
LibriSpeech train-960 corpus, and repeats the PCA plus MiniBatch $K$-means
procedure on encoder hidden states. The iter-2 codebooks use the same
chunk/stride metadata, PCA dimension, and codebook sizes.

\subsection{Self-Supervised Pretraining}
The SSL model is SqueezeFormer-XS with dual prediction heads over the
$K_c=100$ and $K_f=500$ codebooks. During pretraining, spans are sampled
with mask probability 0.30 and mask length 10 target steps. Masked input
frames are replaced by a learnable mask embedding. The loss is the sum of
masked-position cross-entropies for the coarse and fine targets:
\[
  \mathcal{L}_{\mathrm{SSL}} =
  \mathrm{CE}(\hat z^{(c)}, z^{(c)}) +
  \mathrm{CE}(\hat z^{(f)}, z^{(f)}).
\]
Iter-1 is trained for 150k optimizer updates. Iter-2 is initialized from
the iter-1 encoder and trained for 100k additional updates on the refined
targets.

\subsection{CTC Fine-Tuning}
For ASR, the SSL prediction heads are discarded and a CTC projection head
is attached to the SqueezeFormer encoder. The tokenizer is a 128-piece
SentencePiece BPE model trained from LibriSpeech train-960 transcripts by
\code{dataset/train\_tokenizer.py}; the pad token (ID 0) is used as the CTC
blank. The 10h condition uses a reproducible random subset sampled from
\code{train-clean-100} with seed 42, not the standard Libri-Light 10h
split. The 100h condition uses all of \code{train-clean-100}. Test numbers
are greedy CTC decoding on \code{test-clean} and \code{test-other}, with no
external language model.

\section{Hyperparameters}

\begin{table}[h]
\centering
\small
\caption{Feature extraction and target generation settings recovered from
\code{CausalSpecUnit/data.py}, \code{generate\_targets.py}, and the Slurm
target-generation scripts.}
\begin{tabular}{@{}ll@{}}
\toprule
Setting & Value \\
\midrule
Sample rate & 16 kHz \\
Channels & Mono average if needed \\
Waveform normalization & Peak normalization by max absolute value \\
Pre-emphasis & 0.97 \\
Mel bins & 80 \\
FFT size & 512 \\
Window / hop & 25 ms / 10 ms \\
Mel range & 0--8 kHz \\
Log compression & \code{torchaudio.transforms.AmplitudeToDB(top\_db=80)} \\
Normalization & Global CMVN over pretraining corpus \\
Chunk size / stride & 8 / 4 frames \\
PCA dimension & 64, whitened \\
PCA seed & 42 \\
K-means implementation & \code{sklearn.cluster.MiniBatchKMeans} \\
K-means fit sample & Up to 1,000,000 chunks or frames \\
K-means batch size & 8192 \\
K-means max iterations & 300 \\
K-means initialization & \code{n\_init="auto"} \\
Codebooks & $K_c=100$, $K_f=500$ \\
K-means seeds & 42 for coarse, 43 for fine \\
\bottomrule
\end{tabular}
\end{table}

\begin{table}[h]
\centering
\small
\caption{SSL pretraining settings recovered from
\code{slurm/causal\_specunit/02\_pretrain\_ssl\_iter1\_960h\_fair.sh},
\code{04\_iter2\_pretrain\_chain.sh}, and \code{CausalSpecUnit/pretrain\_ssl.py}.}
\begin{tabular}{@{}ll@{}}
\toprule
Setting & Value \\
\midrule
Encoder & SqueezeFormer-XS \\
Encoder parameters & 9,007,782 \todo{confirm final reported rounding} \\
Input dimension & 80 \\
Encoder dimension / layers & 144 / 16 \\
Attention heads & 4 \\
FF / conv expansion & 4 / 2 \\
Conv kernel size & 31 \\
Reduce / recover layers & 7 / 15 \\
Dropout & 0.1 \\
Optimizer & AdamW, betas=(0.9, 0.98), eps=$10^{-9}$ \\
Learning rate & $10^{-3}$ \\
Weight decay & $5\times10^{-4}$ \\
Warmup / hold & 20 epochs / 20 epochs \\
Noam decay rate & 1.0 \\
Batch per process & 128 \\
Distributed processes & 2 \\
Gradient accumulation & 1 \\
Effective batch & 256 utterances \\
Mask probability / span & 0.30 / 10 target steps \\
Gradient clipping & 1.0 \\
Matmul precision & \code{torch.set\_float32\_matmul\_precision("high")} \\
Mixed precision & Not enabled in \code{pretrain\_ssl.py} \\
Iter-1 updates & 150k \\
Iter-2 updates & 100k \\
Training seed policy & \todo{confirm whether model initialization seed is fixed} \\
\bottomrule
\end{tabular}
\end{table}

\begin{table}[h]
\centering
\small
\caption{CTC fine-tuning settings recovered from
\code{slurm/causal\_specunit/06\_test\_table.sh} and
\code{CausalSpecUnit/train\_ctc.py}.}
\begin{tabular}{@{}ll@{}}
\toprule
Setting & Value \\
\midrule
Tokenizer & SentencePiece BPE, 128 pieces \\
CTC blank & Pad token, ID 0 \\
10h split & Custom subset from \code{train-clean-100}, seed 42 \\
100h split & Full \code{train-clean-100} \\
960h split & \todo{confirm exact 960h fine-tuning command/script} \\
10h epochs / batch / accum & 150 / 64 / 1 \\
100h epochs / batch / accum & 100 / 128 / 2 \\
Distributed processes & 2 \\
10h effective batch & 128 utterances \\
100h effective batch & 512 utterances \\
Optimizer & AdamW, betas=(0.9, 0.98), eps=$10^{-9}$ \\
Base / encoder / head LR & $10^{-3}$ / $3\times10^{-4}$ / $10^{-3}$ \\
Weight decay & $5\times10^{-4}$ \\
Warmup / peak epochs & 10 / 50 \\
Noam decay rate & 0.5 \\
Gradient clipping & 1.0 \\
SpecAugment & time=30, freq=20, two masks each \\
SpecAugment schedule & Disabled for final 10 epochs \\
Precision & bfloat16 autocast on CUDA \\
Validation split & \code{dev-other} during training \\
Test decoding & Greedy CTC, no external LM \\
Test sets & \code{test-clean}, \code{test-other} \\
\bottomrule
\end{tabular}
\end{table}

\section{Dataset Split Specification}
The unlabeled pretraining corpus is LibriSpeech train-960:
\code{train-clean-100}, \code{train-clean-360}, and \code{train-other-500}.
The 100h labeled condition is \code{train-clean-100}. The 10h labeled
condition is not a separate official split in this repository; it is
constructed by shuffling utterances from \code{train-clean-100} with
\code{random.Random(42)} and adding utterances until the accumulated audio
duration reaches 10 hours, as implemented by
\code{CausalSpecUnit/data.py::subset\_items\_by\_hours}. The 960h labeled
condition is reported in the main paper, but the exact local launch command
for all 960h rows was not found in the inspected scripts.

\section{Compute Accounting}
The main paper reports 42 GPU-hours for the final two-stage method. The
accounting is:
\[
  \mathrm{GPU\mbox{-}hours} =
  \mathrm{wall\mbox{-}hours} \times \mathrm{number\ of\ GPUs}.
\]
The available scripts use two GPUs for SSL pretraining. The 42 GPU-hour
number is a linear estimate for 250k total SSL updates (150k iter-1 plus
100k iter-2) from a measured 100k-step throughput of about 8.4 wall-hours
on two GPUs. Thus $8.4 \times 2 \times 2.5 \approx 42$ GPU-hours.
The exact GPU model is not recoverable from the local repository:
\todo{confirm GPU model and whether the 8.4h timing came from the final
hardware used for the reported result}. External wav2vec~2.0 and HuBERT
compute values in the main table are published-report figures and are not
matched-hardware measurements.

\section{Additional Results and Robustness}
No multi-seed result table, training-curve export, or completed
phone-purity output file was found in the local repository snapshot. The
main paper should therefore be read as reporting single-run WER values
unless additional logs are added before submission.

\section{Reproducibility Checklist}
\begin{tabular}{@{}p{0.45\linewidth}p{0.45\linewidth}@{}}
\toprule
Item & Status \\
\midrule
Dataset splits defined & Partially; 960h fine-tune launch needs confirmation \\
Preprocessing parameters reported & Yes \\
Target-generation parameters reported & Yes \\
Model architecture specified & Yes \\
SSL optimization specified & Yes, except model seed policy \\
CTC optimization specified & Yes for 10h/100h scripts \\
Decoding specified & Yes, greedy CTC without LM \\
Hardware/accounting specified & Partially; GPU model TODO \\
Seeds reported & Target and 10h subset seeds reported; training seed TODO \\
Code/checkpoint release status & \todo{state anonymous release plan for camera-ready} \\
\bottomrule
\end{tabular}

\end{document}
